\section{Chapter I. Varieties}

\ChapterEntry{err-I-quasiaffine}{\S1. Definition of a quasi-affine variety}{3, 10}

\PassagePage{3}


\Printed{An open subset of an affine variety}
\Revised{A nonempty open subset of an affine variety}

\PassagePage{10}


\Printed{An open subset of a projective variety}
\Revised{A nonempty open subset of a projective variety}

\ChapterEntry{err-I-1-11}{Exercise 1.11. Local complete intersection}{8}


\Printed{which cannot be generated by $2$ elements.}
\Revised{which cannot be generated by $2$ elements even after localization at the origin.}

\ChapterEntry{err-2}{\S2. Test for primality of a homogeneous ideal}{9}


\Printed{homogeneous ideal is prime, it is sufficient to show for any two \emph{homogeneous}}
\Revised{proper homogeneous ideal is prime, it is sufficient to show for any two \emph{homogeneous}}

\ChapterEntry{err-I-2-3}{Corollary 2.3 and Exercise 2.6. Standard affine charts}{11}


\Printed{by the open sets $Y\cap U_i$}
\Revised{by the nonempty open sets $Y\cap U_i$}

\Printed{let $Y_i$ be the affine variety $\varphi_i(Y\cap U_i)$}
\Revised{for each $i$ with $Y\cap U_i\ne\varnothing$, let $Y_i$ be the affine variety $\varphi_i(Y\cap U_i)$}

\ChapterEntry{err-3}{Lemma 3.1. Proof of continuity}{15}


\Printed{A closed set of $\AffSpace[k]{1}$ is a finite set of points,}
\Revised{The inverse image of $\AffSpace[k]{1}$ is all of $Y$, and every proper closed subset of $\AffSpace[k]{1}$ is a finite set of points,}

\ChapterEntry{err-I-3.4-charts}{I, Theorem 3.4. Choice of projective coordinates in the proof}{18}


\Printed{To begin with, let $U_i\subseteq\mathbf P^n$}
\Revised{Choose homogeneous coordinates in which a fixed point of $Y$ has all coordinates nonzero. To begin with, let $U_i\subseteq\mathbf P^n$}

\ChapterEntry{err-I-3.4-localization}{I, Theorem 3.4. Degree-zero parts of localized ideals}{18}


\Printed{$I(Y)S_{(x_i)}$ (cf. Ex.~\Goto{ex-I-2.6}{2.6}), so passing to the quotient, we obtain}
\Revised{$(I(Y)S_{x_i})_0$ (cf. Ex.~\Goto{ex-I-2.6}{2.6}), so passing to the quotient, we obtain}

\Printed{$\varphi_i^*(\mathfrak m'_P)=\mathfrak m_P\cdot S(Y)_{(x_i)}$}
\Revised{$\varphi_i^*(\mathfrak m'_P)=(\mathfrak m_P S(Y)_{x_i})_0$}

\ChapterEntry{err-I-3.19}{I, Exercise 3.19(b). Characteristic and the Jacobian conjecture}{23}


\Printed{The converse of (a) is an unsolved problem, even for $n=2$.}
\Revised{In characteristic zero and for $n=2$, the converse of (a) is an unsolved problem.}

\ChapterEntry{err-I-4.5}{I, Corollary 4.5(ii). Open sets in a birational isomorphism}{26}


\Printed{open subsets $U\subseteq X$ and $V\subseteq Y$ with $U$ isomorphic}
\Revised{nonempty open subsets $U\subseteq X$ and $V\subseteq Y$ with $U$ isomorphic}

\ChapterEntry{err-I-4.9}{I, Proposition 4.9. Irreducible equation of the hypersurface}{27}


\Printed{polynomial equation with coefficients which are rational functions in $x_1,\ldots,x_r$}
\Revised{minimal polynomial equation with coefficients which are rational functions in $x_1,\ldots,x_r$}

\Printed{functions in $x_1,\ldots,x_r$. Clearing denominators, we get an irreducible polynomial $f(x_1,\ldots,x_r,y)=0$.}
\Revised{functions in $x_1,\ldots,x_r$. Clearing denominators and removing the common factor of the coefficients, we get an irreducible polynomial $f(x_1,\ldots,x_r,y)=0$.}

\ChapterEntry{err-I-4-blowup-point}{I, \S4. The birationality assertion for a blow-up at a point}{29}


\Printed{$\varphi$ is a birational morphism of $\widetilde Y$ to $Y$.}
\Revised{if $Y\ne\{O\}$, $\varphi$ is a birational morphism of $\widetilde Y$ to $Y$.}

\ChapterEntry{err-I-4-translation}{I, \S4. Choice of origin in a blow-up}{29}


\Printed{make a linear change of coordinates sending $P$ to $O$.}
\Revised{translate the coordinates to send $P$ to $O$.}

\ChapterEntry{err-I-4.9.1}{I, Example 4.9.1. Characteristic in the two-branch example}{29}


\Printed{Let $Y$ be the plane cubic curve given by the equation $y^2=x^2(x+1)$}
\Revised{Assume that $\operatorname{char} k\ne2$. Let $Y$ be the plane cubic curve given by the equation $y^2=x^2(x+1)$}

\ChapterEntry{err-I-4.4c}{I, Exercise 4.4(c). Singularity type of the cubic}{30}


\Printed{nodal cubic curve $y^2z=x^2(x+z)$ in $\Proj^2$}
\Revised{cubic curve $y^2z=x^2(x+z)$ in $\Proj^2$}

\Printed{$y^2z=x^2(x+z)$ in $\Proj^2$. Show that the projection $\varphi$ from the point $P=(0,0,1)$ to}
\Revised{$y^2z=x^2(x+z)$ in $\Proj^2$ (nodal if $\operatorname{char} k\ne2$). Show that the projection $\varphi$ from the point $P=(0,0,1)$ to}

\ChapterEntry{err-I-4-meet-over-center}{I, \S4. Where distinct tangent directions separate}{30}


\Printed{no longer meet in $X$. Instead, they meet $E$ at points corresponding to the different slopes.}
\Revised{do not meet over $O$. Instead, they meet $E$ at points corresponding to the different slopes.}

\ChapterEntry{p-47}{I, Theorem 5.3. Reference in the proof}{33}


\Printed{5.2}
\Revised{5.2A}

\ChapterEntry{err-I-5-7A}{I.5.7A. Ground field for the common zero}{35}


\Printed{a common zero different from $(0,0,\ldots,0)$}
\Revised{a nonzero common zero over an algebraic closure of $k$}

\ChapterEntry{err-I-5-3-ex}{I, Exercise 5.3. Moving the point to the origin}{36}


\Printed{Make a linear change of coordinates so that $P$ becomes the point $(0,0)$.}
\Revised{Translate the coordinates so that $P$ becomes the point $(0,0)$.}

\ChapterEntry{err-I-5-6a}{I, Exercise 5.6(a). Nonsingularity after the blow-up}{37}


\Printed{the curve $\widetilde Y$ obtained by blowing up $Y$ at $O=(0,0)$ is nonsingular (cf. \Goto{result-I-4.9.1}{(4.9.1)} and Ex.}
\Revised{the curve $\widetilde Y$ obtained by blowing up $Y$ at $O=(0,0)$ is nonsingular at every point over $O$ (cf. \Goto{result-I-4.9.1}{(4.9.1)} and Ex.}

\ChapterEntry{err-I-5-9}{I, Exercise 5.9. Positive degree}{37}


\Printed{Let $f\in k[x,y,z]$ be a homogeneous polynomial, let $Y=Z(f)\subseteq\Proj^2$ be the algebraic set defined by $f$,}
\Revised{Let $f\in k[x,y,z]$ be a homogeneous polynomial of positive degree, let $Y=Z(f)\subseteq\Proj^2$ be the algebraic set defined by $f$,}

\ChapterEntry{err-I-5-11}{I, Exercise 5.11. Elliptic quartic in characteristic two}{38}


\Printed{Let $Y$ be the algebraic set in $\Proj^3$ defined by the equations}
\Revised{Assume $\operatorname{char}k\ne2$. Let $Y$ be the algebraic set in $\Proj^3$ defined by the equations}

\ChapterEntry{err-I-5-12c}{I, Exercise 5.12(c). Singular locus of a quadric}{38}


\Printed{Show that the singular locus $Z=\operatorname{Sing}Q$ of $Q$ is a linear variety (Ex. \Goto{ex-I-2.11}{2.11}) of dimension $n-r-1$.}
\Revised{Show that the singular locus $Z=\operatorname{Sing}Q$ is empty if $r=n$, and otherwise is a linear variety (Ex. \Goto{ex-I-2.11}{2.11}) of dimension $n-r-1$.}

\ChapterEntry{err-I-5-15b}{I, Exercise 5.15(b). Equations for the singular-point locus}{39}


\Printed{degree $d$ in three variables $x,y,z$ has $\binom{d+2}{2}$ coefficients.}
\Revised{positive degree $d$ in three variables $x,y,z$ has $\binom{d+2}{2}$ coefficients.}

\Printed{$\partial f/\partial x_0,\ldots,\partial f/\partial x_n$}
\Revised{$f,\partial f/\partial x,\partial f/\partial y,\partial f/\partial z$}

\ChapterEntry{err-I-6-valuation}{I, \S6. Valuation of zero}{39}


\Printed{$x,y\ne0$, we have}
\Revised{$x,y\ne0$, we have (with the convention $v(0)=+\infty$)}

\ChapterEntry{err-I-6-discrete}{I, \S6. Value group of a discrete valuation}{40}


\Printed{value group $G$ is the integers}
\Revised{value group $v(K\setminus\{0\})$ is order-isomorphic to the integers}

\ChapterEntry{err-I-6-1A-bourbaki}{I, Theorem 6.1A. Reference for the proof}{40}


\Printed{\Goto{bib-474-5}{[2, Ch. VI, \S1, 3]}}
\Revised{\Goto{bib-474-4-item-1}{[1, Ch. VI, \S1, no.~2, Ths.~1--2 and corollary]}}

\ChapterEntry{err-I-6-5}{I.6.5. Localization in the valuation-ring argument}{41}


\Printed{If $R$ is a valuation ring, then $x\notin R$ if and only if}
\Revised{The claim is immediate for $x=0$. Suppose $x\ne0$. If $R$ is a valuation ring, then $x\notin R$ if and only if}

\Printed{Then $\mathfrak n$ is a maximal ideal of $B$, and $B$ is dominated by $R$.}
\Revised{If $\mathfrak n=0$, then every nonzero element of $B$ is a unit in $R$, forcing $R=K$. Thus $\mathfrak n$ is a maximal ideal of $B$, and $B_{\mathfrak n}$ is dominated by $R$.}

\ChapterEntry{err-I-6-nonempty}{I, \S6. Nonempty opens in the abstract-curve construction}{42--43}

\PassagePage{42}


\Printed{$U\subseteq C_K$ is an open subset of $C_K$, we define the}
\Revised{$U\subseteq C_K$ is a nonempty open subset of $C_K$, we define the}

\Printed{Note also by \Goto{result-I-6.5}{(6.5)} that any $f\in K$ is a regular}
\Revised{On the empty set take the zero ring, consisting of the empty function. Note also by \Goto{result-I-6.5}{(6.5)} that any $f\in K$ is a regular}

\Printed{is an open subset $U\subseteq C_K$, where $K$ is a function field of dimension}
\Revised{is a nonempty open subset $U\subseteq C_K$, where $K$ is a function field of dimension}

\PassagePage{43}


\Printed{any open set $V\subseteq Y$, $\mathcal O(V)=\bigcap_{P\in V}\mathcal O_{P,Y}$.}
\Revised{any nonempty open set $V\subseteq Y$, $\mathcal O(V)=\bigcap_{P\in V}\mathcal O_{P,Y}$.}

\ChapterEntry{err-I-6-8}{I.6.8. Domain of regular pullbacks}{44}


\Printed{The uniqueness is clear by construction (it also follows from \Goto{result-I-4.1}{(4.1)}).}
\Revised{\emph{Delete.}}

\Printed{To show that $\overline\varphi$ is a morphism, it will be sufficient to show that regular functions in a neighborhood of $\overline\varphi(P)$ pull back to regular functions on $X$. Let $U_k\subseteq\Proj^n$ be the open set where $x_k\ne0$. Then $\overline\varphi(P)\in U_k$, since $f_{kk}(P)=1$. Now $U_k$ is affine, with affine coordinate ring equal to
\[
k[x_0/x_k,\ldots,x_n/x_k].
\]
These functions pull back to $f_{0k},\ldots,f_{nk}$ which are regular at $P$ by construction. It follows immediately that for any smaller neighborhood $\overline\varphi(P)\in V\subseteq U_k$, regular functions on $V$ pull back to regular functions on $X$. Hence $\overline\varphi$ is a morphism, which completes the proof.}
\Revised{To show that $\overline\varphi$ is a morphism near $P$, take a common open neighborhood $W$ of $P$ on which $f_{0k},\ldots,f_{nk}$ are regular. Since $f_{kk}=1$, these functions give a morphism from $W$ to the affine chart $U_k=\{x_k\ne0\}$, whose coordinate ring is
\[
k[x_0/x_k,\ldots,x_n/x_k].
\]
On $(W\setminus\{P\})\cap\varphi^{-1}(U_j)$, put $a_i=(x_i/x_j)\circ\varphi$. On each such nonempty open set, $a_i=a_k f_{ik}$ as rational functions, hence as regular functions. Taking $i=j$ gives $1=a_k f_{jk}$, so $\varphi$ lies in $U_k$ there and agrees with the constructed morphism. Thus it glues with $\varphi$, with regular pullbacks from $V$ defined on $\overline\varphi^{-1}(V)$. For uniqueness, an extension taking $P$ into $U_j$ makes every $f_{ij}$ regular at $P$, so $r_j$ is minimal. Then $f_{kj}$ is a unit at $P$, and the value in $U_k$ is the one prescribed by the $f_{ik}(P)$.}

\ChapterEntry{err-I-6-12}{I.6.12. Variance of the function-field functor}{45--46}

\PassagePage{45}


\Printed{function fields of dimension $1$ over $k$, and $k$-homomorphisms.}
\Revised{the opposite of the category of function fields of dimension $1$ over $k$ with $k$-homomorphisms.}

\PassagePage{46}


\Printed{clearly inverse to the given functor (i) $\to$ (ii) $\to$ (iii), so we have an}
\Revised{quasi-inverse to the given functor (i) $\to$ (ii) $\to$ (iii), so we have an}

\ChapterEntry{err-I-7-hilbert-module}{\S7. Scope of the Hilbert polynomial}{49, 51}

\PassagePage{49}


\Printed{graded $S$-module, where $S=k[x_0,\ldots,x_n]$.}
\Revised{finitely generated graded $S$-module, where $S=k[x_0,\ldots,x_n]$.}

\PassagePage{51}


\Printed{graded module $M$ over the polynomial ring $S=k[x_0,\ldots,x_n]$.}
\Revised{finitely generated graded module $M$ over the polynomial ring $S=k[x_0,\ldots,x_n]$.}

\ChapterEntry{err-I-7-degree-nonempty}{\S7. Definition of degree}{52}


\Printed{We define the \emph{degree} of $Y$ to be}
\Revised{If $Y\ne\varnothing$, we define the \emph{degree} of $Y$ to be}

\ChapterEntry{err-I-7-7}{Theorem 7.7. Polynomial whose coefficient is compared}{53}


\Printed{$P_i$, we can ignore those $P_i$ of degree $<r-1$.}
\Revised{$P_M$, we can ignore those $P_i$ of degree $<r-1$.}

\ChapterEntry{err-I-7-3}{Exercise 7.3. Tangent line characterized by intersection multiplicity}{54}


\Printed{. We regard the set of lines in $\Proj^2$ as another projective space,}
\Revised{other than a line. We regard the set of lines in $\Proj^2$ as another projective space,}

\ChapterEntry{err-I-8-elliptic}{\S8. Genus one and elliptic curves}{56}


\Printed{Curves with $g=1$ are called \emph{elliptic curves}.}
\Revised{The nonsingular projective representatives of the classes with $g=1$, equipped with a chosen point, are \emph{elliptic curves}.}

\ChapterEntry{err-I-8-genus}{\S8. Arithmetic genus under birational equivalence}{57}


\Printed{in most cases (curves, surfaces, nonsingular varieties in characteristic $0$; see \Goto{result-V-5.6.1}{(V, 5.6.1)}),}
\Revised{among nonsingular projective curves and surfaces, and among nonsingular projective varieties in characteristic $0$ (see \Goto{result-V-5.6.1}{(V, 5.6.1)}),}

\ChapterEntry{err-I-8-picard}{\S8. Divisors and vector bundles}{57}


\Printed{called the \emph{Picard group} of $X$.}
\Revised{which is the \emph{Picard group} of $X$ when $X$ is nonsingular.}

\Printed{algebraic variety. Then one can carry over many constructions from differential geometry}
\Revised{nonsingular algebraic variety. Then one can carry over many constructions from differential geometry}

\ChapterEntry{err-I-8-cohomological-genus}{\S8. Cohomological formula for the genus}{57}


\Printed{curve $X$ can be defined as $\dim H^1(X,\mathcal O_X)$), but it can be used}
\Revised{nonsingular projective curve $X$ can be defined as $\dim H^1(X,\mathcal O_X)$), but it can be used}

\ChapterEntry{err-I-8-residue}{\S8. Residue fields along subvarieties}{58}


\Printed{is not algebraically closed \Goto{ex-I-3.13}{(Ex. 3.13)}, and at times it}
\Revised{need not be algebraically closed \Goto{ex-I-3.13}{(Ex. 3.13)}, and at times it}

